\chapter{CONCLUSION}
\label{ch:conclusion}

\todoresolve{REVIEW SUMMARY OF CONTRIBUTIONS AND FUTURE DIRECTIONS.}

\begin{aigenblock}

\section{Synthesis of Contributions}

Across the six research chapters of this dissertation, we have examined how the representation and use of information govern both the statistical sample complexity of learning from sequential data and the computational complexity of turning those data into inferences and decisions. Rather than treating statistical efficiency and computational tractability as separate problems, each chapter traced how structural choices---how historical data are combined with new interaction, how policies and ensembles are parameterized, and how latent states and future information are approximated---create explicit tradeoffs between sample size, approximation bias, and computation.

\subsection{Using Existing Data to Reduce Exploration Complexity}

Chapter~\ref{ch:hybrid-rl} studied how offline data reduce the statistical burden of online exploration when historical transitions come from the same MDP:
\begin{itemize}
\item \textbf{Direct state--action coverage in a shared environment (Chapter~\ref{ch:hybrid-rl}):} When historical transitions come from the same MDP, their utility is determined by which state--action pairs or feature directions the behavior policy covered. Under general function approximation, incorporating offline transitions into an optimistic Bellman confidence set (Algorithm~\ref{alg:ch02-alg1}) allows the learner to rely on partial offline coverage where it exists and spend online exploration on the remaining Bellman eluder dimension without knowing the coverage boundary in advance. In linear MDPs, targeting poorly covered feature directions during reward-agnostic online exploration before pessimistic offline policy fitting (Algorithm~\ref{alg:ch02-alg2}) augments the offline information matrix from $\Sigma_{\mathrm{off},h}^*$ to $\Sigma_{\mathrm{off},h}^* + \Sigma_{\mathrm{on},h}^*$, while optimistic offline-to-online value iteration (Algorithm~\ref{alg:ch02-alg3}) reduces the online exploration dimension from $d$ to the unspanned subspace dimension $d_{\mathrm{on}}$ in the large-sample regime.
\end{itemize}

\subsection{Representing Policies and Sequential Ensembles Tractably}

Chapters~\ref{ch:actor-critic},~\ref{ch:diffusion-policies}, and~\ref{ch:boosting-inference} examined how the functional representation of a sequentially updated policy or regression ensemble affects both statistical concentration and computational cost:
\begin{itemize}
\item \textbf{Controlling actor policy-class complexity (Chapter~\ref{ch:actor-critic}):} In actor--critic reinforcement learning, unrolling natural policy updates accumulates all past critics into the policy representation, causing the policy class covering number to grow with the number of critic updates and destabilizing policy-evaluation Bellman targets. Targeting the optimal value function $Q^*$ with a rare-switching optimistic critic and resetting the actor to a uniform distribution at each switch (Algorithm~\ref{alg:NORA}) decouples the Bellman target operator from actor drift and achieves optimal $\widetilde{O}(\sqrt{T})$ regret and $\widetilde{O}(\epsilon^{-2})$ sample complexity under general function approximation, while hybrid extensions (Algorithms~\ref{alg:NOAH-pi}--\ref{alg:NORA-long}) leverage offline data to either eliminate explicit optimism or sharpen online regret.
\item \textbf{Optimizing continuous multimodal policies via diffusion flows (Chapter~\ref{ch:diffusion-policies}):} When actions are continuous and multimodal, diffusion policies represent complex action distributions through $K$-step denoising paths whose marginal action density is intractable. We showed how vanilla diffusion policy gradients (Algorithm~\ref{alg:dppo-vanilla-pg}), deterministic-flow DDIM natural policy gradients under reset, mixing, or hybrid logged-trajectory access (Algorithm~\ref{alg:dppo-reinforce}), and denoising score-matching behavior cloning of an optimistic or hybrid critic (Algorithm~\ref{alg:dipo-bc}) each trade off sampling access assumptions against comparison-policy, advantage-transfer, score-matching, and sampler errors.
\item \textbf{Asymptotic normality in sequentially boosted ensembles (Chapter~\ref{ch:boosting-inference}):} In supervised gradient boosting and explainable boosting machines, sequential residual fitting couples every tree to all previous trees. Combining moving-average shrinkage with random tree dropout (Algorithm~\ref{alg:random-dropout}), parallel structured dropout (Algorithm~\ref{alg:structured-dropout}), or cyclic univariate feature updates (Algorithm~\ref{alg:brebm_ident}) yields contraction mappings whose fixed points suffer less shrinkage bias toward zero while admitting explicit central limit theorems for confidence intervals, prediction intervals, and hypothesis tests.
\end{itemize}

\subsection{Inference and Control under Partial Observability}

Chapters~\ref{ch:pomp-estimation} and~\ref{ch:pomdp-control} relaxed the assumption that the learner observes the underlying state directly, making sequential Monte Carlo filtering part of the estimation and control loop:
\begin{itemize}
\item \textbf{Plug-and-play differentiable particle filtering for parameter estimation (Chapter~\ref{ch:pomp-estimation}):} To resolve the discontinuity of resampled particle filters without sacrificing the plug-and-play property or suffering linear-in-$N$ particle weight-variance growth, our fixed-ancestor weight-discounted particle filter (Algorithm~\ref{alg:mop}) and its single-pass differentiable counterpart (Algorithm~\ref{alg:dmop}) introduce a discount factor $\alpha \in [0,1]$ on carried log-weight derivatives, achieving $O(N\alpha^{2/(1-\alpha)})$ squared score bias and $O(N(1+\log J)^2/(J(1-\alpha)^2))$ particle variance under uniform mixing. Pairing iterated filtering warm starts with discounted automatic-differentiation optimization (Algorithm~\ref{alg:ifad}) combines global mode discovery with rapid local gradient refinement on high-dimensional epidemiological and ecological dynamical systems.
\item \textbf{Decomposing estimation, filtering, rollout, and lookahead errors in POMDP control (Chapter~\ref{ch:pomdp-control}):} For online decision making in POMDPs, Algorithm~\ref{alg:pf_mcts} integrates online trajectory maximum likelihood estimation, $J$-particle bootstrap filtering, and depth-$d$ Monte Carlo belief-tree planning (Algorithm~\ref{alg:mctq}) with $N$ rollouts and a state-value QMDP leaf bootstrap. The resulting regret bound isolates the statistical model-learning rate $\widetilde{O}(T^{-1/4})$ from the particle filtering error $O(J^{-1/(2\kappa_{\mathrm{pf}})})$, rollout estimation error $O(C_{\mathrm{mc}}(d,M) N^{-1/2})$, and structural lookahead error $O(\gamma^{d+1}\overline{\Xi}_T)$, clarifying when deeper belief-tree lookahead is needed to capture the value of information-gathering actions and when a depth-zero particle state-rollout approximation suffices.
\end{itemize}

\section{Future Directions}

The results in this dissertation point toward several concrete open problems at the interface of statistical learning theory, stochastic filtering, and scalable sequential decision making:
\begin{enumerate}
\item \textbf{Closing sample-size threshold and horizon gaps in hybrid and actor--critic RL:} In Chapter~\ref{ch:hybrid-rl}, the offline-to-online regret guarantee for Algorithm~\ref{alg:ch02-alg3} in linear MDPs requires a large-sample burn-in threshold of order $\widetilde{\Omega}(d^{13} H^{14})$ to align the empirical variance-weighted covariances with their population counterparts. Similarly, in Chapter~\ref{ch:actor-critic}, while Algorithm~\ref{alg:NORA} achieves the optimal $\epsilon^{-2}$ accuracy exponent under general function approximation, its horizon dependence remains $\widetilde{O}(d H^5 \log|\mathcal{A}| \epsilon^{-2} + d H^4 \log|\mathcal{F}| \epsilon^{-2})$. Developing variance-aware actor--critic and hybrid updates that achieve minimax horizon dependence without large burn-in thresholds remains an important theoretical question.
\item \textbf{Adaptive tree structure and growing-depth inference in boosted ensembles:} The central limit theorems in Chapter~\ref{ch:boosting-inference} rely on finite-partition regularity conditions to establish contraction of the expected kernel operator as the number of boosting rounds grows. Extending valid frequentist confidence intervals to adaptive splitting rules with growing tree depth and honesty across ancestral sample reuse is a natural next step for nonparametric ensemble inference.
\item \textbf{Joint gradient-based estimation and control in high-dimensional partially observed systems:} Chapter~\ref{ch:pomp-estimation} developed low-MSE pathwise score estimators for passive partially observed Markov processes using discounted automatic differentiation, while Chapter~\ref{ch:pomdp-control} analyzed online POMDP control using trajectory maximum likelihood estimation and bootstrap particle belief planning. Combining discounted differentiable particle filtering directly with online policy optimization and adaptive belief-tree allocation---particularly in high-dimensional spatiotemporal systems where localized or block particle resampling is needed to prevent weight collapse---promises to further unify statistical inference and sequential control under partial observability.
\end{enumerate}
\end{aigenblock}
